\documentclass[11pt]{article}
\usepackage[T1]{fontenc}
\usepackage{textcomp}
\usepackage[margin=0.75in]{geometry}
\usepackage{amsmath,amssymb,amsthm}
\usepackage{array,booktabs}
\usepackage{hyperref}
\setlength{\parindent}{0pt}
\newtheorem{theorem}{Theorem}
\theoremstyle{definition}
\newtheorem{definition}{Definition}
\title{Flow Proof Artifact: prop-08-fourfold-rectangle-plus-square}
\author{Generated by \texttt{flow doc proof}}
\date{Flow Proof Book}
\begin{document}
\maketitle
If a straight line is cut at random, four times the rectangle by the whole and one segment together with the square on the remaining segment equals the square on the whole and that segment as on one straight line.\\[0.5em]
\textbf{Source.} euclid — Elements, Book II, Proposition 8\\[0.5em]
\bigskip
\noindent{\large \textbf{Derived fact 1.} \textit{If a straight line is cut at random, four times the rectangle by the whole and one segment together with the square on the remaining segment equals the square on the whole and that segment as on one straight line} \hfill the Euclidean plane \cdot Euclid Book II \cdot Proposition 8: four times a rectangle plus the square on the remainder equals a square on one line}\par
\smallskip\noindent\textit{Coordinate: the Euclidean plane \cdot Euclid Book II \cdot Proposition 8: four times a rectangle plus the square on the remainder equals a square on one line}\par
\smallskip\noindent\textit{Source: euclid — Elements, Book II, Proposition 8}\par
\smallskip\noindent\textit{Built on: proposition 7: the square on the whole plus that on one segment equals twice a rectangle plus a square, for Euclid Book II on the Euclidean plane}\par
\medskip
\noindent\textbf{Goal.} If a straight line is cut at random, four times the rectangle by the whole and one segment together with the square on the remaining segment equals the square on the whole and that segment as on one straight line\par
\noindent$\displaystyle 4 \cdot rect(whole, C) + sq(D) = sq(\text{whole and C})$\par
\medskip
\noindent
\begin{tabular}{@{} >{\raggedright\arraybackslash}p{0.06\textwidth} >{\raggedright\arraybackslash}p{0.47\textwidth} >{\raggedleft\arraybackslash}p{0.41\textwidth} @{}}
\textbf{\#} & \textbf{Proof} & \textbf{Mathematics} \\
\midrule
\textbf{1.} & We prove that proposition 8: four times a rectangle plus the square on the remainder equals a square on one line for Euclid Book II the Euclidean plane. & \\[0.45em]
\textbf{2.} & Let whole = straight\_line\_AB. & \\[0.45em]
\textbf{3.} & Let C = segment\_AC. & \\[0.45em]
\textbf{4.} & Let D = segment\_CB. & \\[0.45em]
\textbf{5.} & From step 2, step 3, and step 4, this implies 4 times rect(whole, C) plus sq(D) equals sq(the conjunction of whole and C). Hence proven. & $\displaystyle 4 \cdot rect(whole, C) + sq(D) = sq(\text{whole and C})$ \\[0.45em]
\textbf{6.} & We invoke the derived fact governing Euclid Book II the Euclidean plane: proposition 7: the square on the whole plus that on one segment equals twice a rectangle plus a square, for Euclid Book II on the Euclidean plane. & \\[0.45em]
\bottomrule
\end{tabular}
\label{thm:Geometry-Euclid Book II-proposition_8_four_times_a_rectangle_plus_the_square_on_the_remainder_equals_a_square_on_one_line}
\par\medskip\hrule
\end{document}
